\section{Methods}
\label{sec:methods}

\subsection{Overview}
We evaluate whether music foundation models align with human auditory representations and whether this alignment relates to downstream music task performance. The pipeline has three stages: (1) estimate brain alignment using natural-sound fMRI datasets, (2) extract layer-wise model embeddings and probe them on music tasks, and (3) quantify the correspondence between brain alignment and downstream task performance. The implementation follows the auditory alignment protocol in Tuckute et al. (2023) and uses the downstream task taxonomy summarized in Yuan et al. (ISMIR, 2025/2026), with consistent preprocessing and evaluation across models.

\subsection{Datasets}
\paragraph{Natural sound fMRI.}
We use two publicly available natural sound datasets that have been used in prior auditory alignment work: Norman-Haignere et al. (2015; NH2015) and Boebinger et al. (2021; B2021). Both provide stimulus-level fMRI responses to a set of 165 natural sounds with three repetitions per stimulus. We follow the original train/test splits and reliability estimates provided with the datasets and reproduce the preprocessing and evaluation protocol from Tuckute et al. (2023).

\paragraph{Music downstream tasks.}
We evaluate model embeddings on a set of music tasks spanning tagging, emotion, and genre/instrument recognition. The current task set includes:
\begin{itemize}
  \item EMO: music emotion regression with separate valence and arousal targets.
  \item MTG Jamendo: Top50 tagging, Mood/Theme, Genre, and Instrument subsets.
\end{itemize}
All datasets are stored under \texttt{./dataset/music\_dataset} with manifest files defining splits and label targets. For each task, we use the official or widely adopted train/validation/test splits from dataset providers.

\subsection{Models}
We evaluate a set of pretrained audio and music foundation models. The current suite includes:
\begin{itemize}
  \item CLAP (general audio) and CLAP variants (music-only, music+speech).
  \item Mulan.
  \item MERT (95M and 330M).
  \item MusicGen (small, medium, large, and large-melody).
\end{itemize}
All models are used in frozen mode without finetuning. For each model, we extract representations from every encoder block (layer-wise) and the final pooled embedding. Model checkpoints are obtained via Hugging Face, and all embedding extraction is GPU-accelerated.

\subsection{Embedding Extraction}
Audio is loaded as mono and resampled to each model's native sampling rate using polyphase resampling. For each clip, we run the corresponding encoder and export both the final embedding and layer-wise embeddings. We apply model-specific pooling to convert hidden states to one vector per layer, and save embeddings to disk with metadata including model name, layer index, dataset, task, and split. Extraction is resume-safe and does not recompute completed partitions.

\paragraph{CLAP.}
We use the CLAP audio encoder (HTSAT) and compute the standard audio embedding by applying the audio projection head to the pooled audio representation and L2-normalizing the result. For layer-wise analysis, we extract hidden states from the CLAP audio model and average across time or time-frequency dimensions, matching the encoder's default pooling behavior. We evaluate the general CLAP model as well as music-only and music+speech variants.

\paragraph{MuLan.}
We use the MuQ-MuLan audio encoder at 24 kHz. For the final embedding, we obtain audio latents for all clips of a recording and average across clips. For layer-wise analysis, we extract backbone hidden states and average across time.

\paragraph{MERT.}
We use MERT (95M and 330M) with its Wav2Vec2-style feature extractor. We compute pooled representations by attention-masked mean pooling over time. Layer-wise embeddings are computed by pooling each hidden layer with the same attention mask.

\paragraph{MusicGen.}
We extract decoder-side hidden states from MusicGen by encoding audio into EnCodec latents and running the decoder on these latents with unconditional text inputs. Prompt chunking is disabled to preserve a single frame of audio codes. We pool decoder hidden states by averaging across the sequence dimension. We evaluate small, medium, large, and melody-large variants.

\subsection{Brain Alignment Analyses}
\paragraph{Linear encoding (regression).}
We predict voxel responses from model embeddings using ridge regression with cross-validation, following Tuckute et al. (2023). Let $X \in \mathbb{R}^{N \times D}$ be model embeddings for $N=165$ stimuli and $Y \in \mathbb{R}^{N \times V \times 3}$ be voxel responses with three repetitions. For each of 10 random splits (seed 0), we sample 83 train and 82 test stimuli without replacement. We fit \texttt{RidgeCV} with an alpha grid $\{10^{-50},\dots,10^{49}\}$ and \texttt{alpha\_per\_target=True}, and predict voxel responses from the mean of the three repetitions. We compute $R^2$ as the squared Pearson correlation between predictions and observed responses, after clipping correlations to be non-negative. 

Noise-corrected $R^2$ follows the same correction used in Tuckute et al.: we compute pairwise split-half correlations across repetitions for observed responses ($r_v$) and predicted responses ($\hat{r}_v$), apply the Spearman-Brown adjustment $r \leftarrow \frac{3r}{1+2r}$, and floor $r_v \ge 0.182$ and $\hat{r}_v \ge 0.183$. The corrected score is
\[
R^2_{\mathrm{corr}} = \frac{R^2}{r_v \cdot \hat{r}_v},
\]
clipped to be non-negative. We report the median across splits for each voxel, and the median across voxels for dataset-level summaries. Subject-level summaries are computed by randomly masking half of the 10 splits per voxel, taking the median across the remaining splits, aggregating by subject, and averaging across 10 iterations. The spectrotemporal filter model from Tuckute et al. is included as a baseline; its performance is shown as a dashed reference line in figures.

\paragraph{Representational similarity analysis (RSA).}
We compute representational dissimilarity matrices (RDMs) for each model layer and for human fMRI responses. For each split (same 10 splits as regression), we z-score features across stimuli (mean subtract and divide by standard deviation) and compute Pearson correlation matrices. RDMs are defined as $1 - \rho$, where $\rho$ is the Pearson correlation between stimulus representations. RSA scores are Spearman correlations between the upper-triangular entries of model and human RDMs. 

We compute RSA per participant by comparing model test RDMs to participant test RDMs; participant scores are the median across splits. We report the median Spearman $\rho$ across participants as the RSA summary. We also compute leave-one-out neural reliability by correlating each participant's RDM with the mean RDM of the remaining participants. Reported values and plots follow Tuckute et al. (2023), including the same spectrotemporal baseline reference.

\subsection{Downstream Probing}
We evaluate embeddings using linear probes with fixed model representations. For each dataset, embeddings are standardized with \texttt{StandardScaler} before probing. We do not tune hyperparameters; all probes use $\alpha=1.0$.

\paragraph{Single-label classification.}
We train a \texttt{RidgeClassifier} on train+valid and evaluate on test. We report accuracy (primary) and macro-F1.

\paragraph{Multi-label tagging.}
We train a one-vs-rest \texttt{RidgeClassifier} with decision function scores. Predictions are obtained by thresholding at zero. We report macro average precision (primary), macro ROC-AUC, and macro-F1.

\paragraph{Regression.}
We train a \texttt{Ridge} regressor on train+valid and evaluate on test. For multi-target regression (EMO), valence and arousal are treated as separate targets. We report Pearson and Spearman correlations, $R^2$, and MAE for each target, plus their means across targets. The primary metric is the mean Pearson correlation (or rounded accuracy for integer targets).

\paragraph{Evaluation protocol.}
If a dataset provides explicit train/valid/test splits, we use that protocol. For datasets with fold-based splits, we perform circular cross-validation by rotating test and validation folds and report the mean across folds. We report both final-layer and best-layer results, and we provide layer-wise performance curves.

\subsection{Brain--Behavior Correlation}
To study the relationship between brain alignment and downstream utility, we compute correlations across models between (i) brain alignment scores (regression and RSA) and (ii) downstream task performance. Correlations are computed separately for each dataset (NH2015, B2021) and for each probing task. We report both Pearson $r$ and Spearman $\rho$, and visualize per-model scatter plots for final-layer and best-layer settings.

\subsection{Implementation Details}
All experiments are run in the \texttt{universal} conda environment on NVIDIA GPUs. The analysis code is organized into three packages under \texttt{./ismir2026}: \texttt{natural\_sound\_align} (brain alignment), \texttt{music\_task\_embed} (embedding extraction), and \texttt{music\_task\_probe} (probing and correlation analysis). Random splits use a fixed seed (0) for reproducibility. The pipeline is resume-safe and deterministic given fixed seeds, enabling incremental expansion to additional models and tasks.
